\documentclass[11pt]{article}
\usepackage[T1]{fontenc}
\usepackage[utf8]{inputenc}
\usepackage[english]{babel}
\usepackage[margin=1in]{geometry}
\usepackage{amsmath,amssymb}
\usepackage{enumitem}
\usepackage{microtype}

\ifdefined\pdfcatalog
  \pdfcatalog{/Lang (en-US)}
\fi

\setlength{\parindent}{0pt}
\setlength{\parskip}{0.55em}
\setlength{\emergencystretch}{2em}
\setlist{nosep,leftmargin=*}

\newcommand{\findinglabel}[1]{\textbf{#1}}
\newcommand{\classification}[1]{\textit{Classification: #1.}}

\title{Technical Review of ``Early-Stopping Activation Steering for Factual Preference Alignment in Vietnamese Medical Language Models''}
\author{Manuscript review}
\date{August 25, 2026}

\begin{document}
\selectlanguage{english}
\maketitle

\section*{Overall assessment}

The manuscript has improved the natural-completion narrative and now explicitly describes its metric trade-off.  It also expands the BM25 retrieval report with Recall@$k$, MRR, an oracle-context BERTScore, and updated activation values.  The bounded 200-token McNemar table remains internally consistent and continues to be the strongest supported result.

The new revision is still not submission-ready because several newly updated headline numbers are arithmetically or internally wrong.  The activation section claims distances that do not follow from its own norms, the BM25 latency percentage and absolute delay do not match 7.32~s versus 14.45~s, and the asserted $2.67\times$ context expansion contradicts the table's 75.4 versus 118.0 prompt tokens.  The corpus size also changes from 4,495 to 14,576 passages, and the Discussion retains obsolete natural-completion values.  More fundamentally, norm elevation is still treated as saturation, RefPref is still interpreted as clinical factuality, and the placebo, latency, dataset, vector, and implementation analyses remain incomplete.  Major revision is required.

\section*{Major technical findings}

\subsection*{1. The activation-diagnostic arithmetic is internally wrong}

\classification{Definite quantitative errors}

\findinglabel{Location.} Abstract; contribution~3; Empirical Representation Saturation Diagnostics.

\findinglabel{Problem.} At $t=50$, the manuscript reports baseline norm 54.21 and continuous-steering norm 56.18, but calls the difference $+2.93$.  The actual difference is
\[
56.18-54.21=1.97.
\]
It also reports linear-decay norm 51.20 and states that it is ``within 0.19 units'' of the 54.21 baseline.  The actual signed and absolute differences are
\[
51.20-54.21=-3.01,
\qquad
|51.20-54.21|=3.01.
\]
Thus linear decay does not converge to the displayed baseline at $t=50$; it undershoots it substantially.  The Abstract nevertheless says it ``smoothly converges back to baseline levels.''  The hard-cutoff claim gives a change from 56.28 to 52.75 without identifying the two decoding steps.

\findinglabel{Why it matters.} These values are presented as direct empirical support for the central mechanism.  Incorrect subtraction reverses the interpretation of the linear-decay trajectory and invalidates the headline stabilization claim.

\findinglabel{Suggested fix.} Regenerate all activation summaries directly from the recorded trajectories and automatically test every displayed difference.  Report the time index, mean, dispersion, confidence interval, and per-step sample count for each condition.  Replace the present stabilization conclusion unless corrected data demonstrate convergence.

\subsection*{2. Activation norm elevation does not establish saturation or manifold departure}

\classification{Unsupported mechanism claim and implementation ambiguity}

\findinglabel{Location.} Abstract; Introduction; Early-Stopping Steering Hook \& Saturation Diagnostics; Empirical Representation Saturation Diagnostics.

\findinglabel{Problem.} Even after correcting the arithmetic, a larger scalar $\ell_2$ norm is not equivalent to representation saturation.  Saturation requires an operational criterion such as clipping, nonlinear loss of responsiveness, directional collapse, or a measured degradation caused by the altered state.  Likewise, a scalar radius cannot show that a point is outside a high-dimensional ``activation manifold.''  Because early steering is explicitly set to zero after $K=16$, removal of the direct perturbation at later steps is expected by construction and is not itself evidence that decay repairs a saturation process.

The paper does not state whether it measures pre-hook $h_l^{(t)}$ or post-hook $\hat h_l^{(t)}$ consistently, whether it averages norms or takes the norm of an average, which tensor positions are modified, how prompts ending before $t=50$ or $t=100$ are handled, or whether conditions are compared on identical token contexts.  Free-generation trajectories are confounded because different schedules produce different previous tokens.

\findinglabel{Why it matters.} The central novelty is motivated as prevention of a specific internal failure, but the diagnostic presently measures only an under-specified norm difference.

\findinglabel{Suggested fix.} Define saturation before analysis.  Report complete pre/post-hook trajectories with uncertainty, survivor counts, projection on $v_{\text{steer}}$, cosine drift, covariance or distributional distance, and response to an additional controlled perturbation.  Add a fixed-token replay or teacher-forced analysis so conditions share the same context.  Link activation diagnostics to repetition and factual errors.  Use ``norm elevation'' rather than ``saturation'' or ``outside the manifold'' unless those stronger properties are directly demonstrated.

\subsection*{3. The RAG latency and context-overhead claims contradict the displayed measurements}

\classification{Definite arithmetic and data-definition errors}

\findinglabel{Location.} Abstract; BM25 RAG Results; Table~\texttt{tab:baselines}; Conclusion.

\findinglabel{Problem.} The manuscript compares 7.32~s with 14.45~s but reports an 87.1\% increase and a 6.38~s delay.  The correct values from the displayed means are
\[
14.45-7.32=7.13\ \text{s},
\qquad
\left(\frac{14.45}{7.32}-1\right)100\%\approx97.4\%.
\]
The Abstract also retains a $2.67\times$ prompt-context expansion.  Table~\texttt{tab:baselines}, however, gives 118.0 prompt tokens for baseline and only 75.4 for BM25 RAG:
\[
\frac{75.4}{118.0}\approx0.639,
\]
which is a 36.1\% reduction, not an expansion.  If 75.4 is only the retrieved passage length rather than the full prompt length, the table label and comparison are wrong.  The Conclusion repeats the obsolete 6.38~s value and calls 75.4 tokens an advantage comparison without resolving the denominator.

\findinglabel{Why it matters.} Efficiency is a headline contribution.  The displayed numbers cannot support the stated latency and memory conclusions until timing and token-count definitions are reconciled.

\findinglabel{Suggested fix.} Recompute the systems table from raw logs and define whether prompt tokens include instructions, user query, and retrieved context.  Report retrieval, tokenization, prefill, and decoding time separately, with repeated-run uncertainty.  Generate all percentages and deltas automatically from the table source.  Remove the $2.67\times$, 87.1\%, and 6.38~s claims unless the underlying measurements are changed and documented.

\subsection*{4. The BM25 experiment remains internally inconsistent and insufficient for a general RAG comparison}

\classification{Definite reporting inconsistency and likely weak baseline}

\findinglabel{Location.} Abstract; Experimental Setup; BM25 RAG Results; Table~\texttt{tab:baselines}.

\findinglabel{Problem.} Experimental Setup defines retrieval over 4,495 passages, whereas Results uses 14,576 passages.  The Abstract describes ``19.20\% Recall@1 retrieval noise,'' although Recall@1 is a success rate; the corresponding miss rate is 80.80\%.  The manuscript gives Recall@1/3/5 and MRR but still omits passage construction, Vietnamese tokenization, query text, BM25 implementation and parameters, relevance mapping, and retrieval-failure policy.  Oracle Gold-Context RAG appears only as BERTScore 0.8002, with no RefPref, latency, prompt length, generation settings, or paired analysis.  No paired interval/test compares BM25 RAG with baseline or steering.

\findinglabel{Why it matters.} The corpus discrepancy prevents reproduction, and a low-recall, under-specified configuration cannot justify a general claim that steering is superior to RAG.

\findinglabel{Suggested fix.} Use one documented corpus and publish its passage IDs and construction rules.  Report miss rate separately from recall, retrieval metrics for each category, generation results conditional on retrieval success, complete oracle results, and stronger lexical or dense retrievers.  Compare methods on identical records using paired uncertainty.  Describe the current result as one BM25 configuration rather than a general RAG baseline.

\subsection*{5. The natural-completion summary is improved, but the Discussion remains obsolete and no paired uncertainty is reported}

\classification{Definite section-to-section contradiction and statistical omission}

\findinglabel{Location.} Abstract; Primary Natural Completion Evaluation; Table~\texttt{tab:long\_gen}; Discussion; Conclusion.

\findinglabel{Problem.} The Results and Conclusion now consistently report baseline 65.40\%, linear decay 67.80\%, continuous steering 68.60\%, and hard cutoff 70.60\%.  Discussion still states a ``positive but non-statistically confirmed preference trend (+1.00~pp),'' which belongs to an older experiment.  It also says natural completion achieves a 100.00\% EOS rate, whereas continuous steering and linear decay are 99.80\%.  The paragraph before the table similarly says both baseline and early-stop models achieve 100.00\%, contradicting the next paragraph.  No paired contingency tables, confidence intervals, or multiplicity policy are provided for the three schedule-versus-baseline comparisons.

\findinglabel{Why it matters.} The primary evaluation cannot support significance, schedule ranking, or robustness claims from marginal percentages alone, and stale Discussion text changes the reported effect size.

\findinglabel{Suggested fix.} Replace the obsolete Discussion and EOS statements.  Report raw counts, full paired transitions, exact tests, paired CIs for RefPref/BERTScore/Rep-4, and a prespecified multiplicity approach.  State that 99.80\% corresponds to one non-EOS output out of 500 and report its finish reason.

\subsection*{6. Natural completion shows a substantial metric trade-off, not uniform quality improvement}

\classification{Unsupported comparative language}

\findinglabel{Location.} Abstract; Primary Natural Completion Evaluation; Discussion; Conclusion.

\findinglabel{Problem.} Every steering schedule increases RefPref but lowers reference BERTScore and increases Rep-4 relative to baseline.  For hard cutoff, Rep-4 rises from 4.10\% to 5.35\%, an absolute increase of 1.25 points and a relative increase of approximately 30.5\%.  For linear decay, the relative increase is approximately 31.0\%.  Describing these as ``minor'' increases is not justified without uncertainty or a practical threshold.  Hard cutoff has the best RefPref, baseline has the best BERTScore and Rep-4, and linear decay does not dominate continuous steering.  The Conclusion nevertheless says linear decay yields consistent improvements.

\findinglabel{Why it matters.} A reader could interpret ``improvement'' as general output quality or clinical correctness, although the result is metric-dependent and early stopping increases the paper's own repetition measure.

\findinglabel{Suggested fix.} Present all paired metric differences with CIs and explicitly frame the result as a trade-off.  Define practical noninferiority margins for BERTScore and Rep-4 if the intended claim is that RefPref improves without unacceptable degradation.  Select the primary schedule on validation data and avoid ``consistent improvement'' unless it applies to a named endpoint.

\subsection*{7. RefPref remains a semantic proxy rather than clinical factual correctness}

\classification{Unsupported interpretation of a disclosed proxy}

\findinglabel{Location.} Abstract; Introduction; contributions; Results; Conclusion; keywords.

\findinglabel{Problem.} Methods correctly discloses that RefPref is not clinician-adjudicated correctness, but the manuscript still calls the vector ``truthfulness,'' describes RefPref as factual alignment and accuracy, and retains hallucination-mitigation language.  A response can be closer to $y^+$ than one synthetic $y^-$ while containing a wrong dose, unit, negation, contraindication, or interaction mechanism.  The natural experiment's simultaneous RefPref increase and reference-BERTScore decrease further demonstrates that the endpoint is not a direct correctness measure.

\findinglabel{Why it matters.} Clinical safety and deployability cannot be inferred from relative embedding similarity against two templates.

\findinglabel{Suggested fix.} Add blinded clinician-adjudicated factual correctness and harmful-error rates, stratified by category.  Validate RefPref against those labels with sensitivity, specificity, agreement, and failure cases.  Until then, use ``reference-preference rate'' consistently and avoid clinical-accuracy or hallucination-mitigation claims as measured outcomes.

\subsection*{8. The random-direction analysis remains statistically overstated}

\classification{Unsupported strength of evidence}

\findinglabel{Location.} Abstract; contribution~4; directional-control Methods; placebo Results; Conclusion.

\findinglabel{Problem.} The standardized distance $(77.20-56.50)/5.11\approx4.05$ is not a calibrated $4\sigma$ hypothesis test.  With 20 sampled directions and none reported above 63.00\%, the simplest corrected empirical randomization bound is $1/(20+1)\approx0.0476$, assuming a prespecified exchangeable procedure.  Seeds, per-direction results, BERTScore dispersion, and hierarchical uncertainty over questions and directions are not reported.  Negative steering produces a $-10.20$-point change while positive steering produces $+4.20$ points, so the effect is not symmetric despite the stated control purpose.

\findinglabel{Why it matters.} The controls provide direction-specific evidence but do not prove a factual representation or support a Gaussian-tail interpretation.

\findinglabel{Suggested fix.} Release seeds and per-direction/per-question outputs.  Report mean, sample SD, range, and CIs for every metric, use an empirical randomization test and hierarchical bootstrap, and analyze the positive/negative asymmetry.  Replace the $4\sigma$ language with the actual test result.

\subsection*{9. Latency, memory, ablation, and matched-dose claims remain under-specified}

\classification{Systems ambiguity and formula inconsistency}

\findinglabel{Location.} cumulative-dose equations; Experimental Setup; Tables~\texttt{tab:baselines} and \texttt{tab:equal\_alpha\_ablation}; Discussion.

\findinglabel{Problem.} Every non-RAG condition is still reported as exactly 7.32~s, without trial count, dispersion, warm-up, synchronization, batch size, GPU count, timing boundary, or peak-memory measurement.  ``No measurable overhead'' is not supported by an equivalence interval.  The ablation emphasizes differences of 0.0005--0.0018 BERTScore without paired CIs or multiplicity handling, while $D_2$ and temporal placement remain unmatched.  Finally, $D_{\text{decay}}$ is an absolute nonnegative dose, but
\[
\alpha_{\text{match}}=D_{\text{decay}}/T
=\frac{K+1}{2T}\alpha_0
\]
mixes an absolute quantity with a signed one when $\alpha_0<0$.

\findinglabel{Why it matters.} The efficiency claim is not statistically measured, small schedule rankings may be noise, and the matched-dose definition is not valid for the negative-steering domain stated elsewhere.

\findinglabel{Suggested fix.} Run repeated component-wise timing and memory measurements with uncertainty and an equivalence margin.  Report paired ablation CIs with an exploratory multiplicity policy.  Define $\alpha_{\text{match}}=\operatorname{sgn}(\alpha_0)D_1/T$ with $\sum_t|\alpha_{\text{match}}|=D_1$, or explicitly restrict the definition to $\alpha_0>0$.

\subsection*{10. Dataset, vector, layer, metric, and hook specifications remain incomplete}

\classification{Reproducibility limitation and likely leakage risk}

\findinglabel{Location.} benchmark construction; vector estimation; layer sweep; metric definitions; Experimental Setup.

\findinglabel{Problem.} ``Expert-structured'' remains undefined, and the manuscript provides no annotation protocol, counter-answer construction rules, quality checks, deduplication, source IDs, licensing, or source/template-grouped splits.  The 500-item sample rule and seed are absent.  The direction is extracted at the final token of a completed answer but injected into early-generation states, without testing prompt-only, initial-token, or pooled vectors.  The layer experiment remains a steering sweep reported only as a range, yet it is called subspace probing and used to infer distributed truthfulness.  Rep-4, Vietnamese ROUGE tokenization, full BERTScore settings, statistical software calls, hook module/indexing, modified tensor positions, KV-cache behavior, versions, seeds, and GPU count remain unspecified.

\findinglabel{Why it matters.} Template/source leakage and implementation choices could affect every reported result and prevent independent reproduction.

\findinglabel{Suggested fix.} Add a dataset card and grouped split manifests, publish the sample seed, document expert validation, test alternative vector positions and layers, provide complete metric/statistical configurations, and release executable hook pseudocode plus per-question outputs.

\section*{Equation, notation, sign, branch, and dimensional audit}

The direction $\mu_l^+-\mu_l^-$ and the intervention $+\alpha(t)v_{\text{steer}}^{(l)}$ are sign-consistent.  Unit normalization makes the vector dimensionally compatible with the hidden state when $\alpha(t)$ is an activation-scale coefficient.  The linear schedule and positive-$\alpha_0$ values of $D_{\text{continuous}}$, $D_{\text{cutoff}}$, and $D_{\text{decay}}=(K+1)|\alpha_0|/2$ are arithmetically correct.  No inverse-trigonometric functions, coordinate transformations, or branch/quadrant choices occur.

The concrete issues are the signed/absolute mismatch in $\alpha_{\text{match}}$, ambiguity between $h_l^{(t)}$ and $\hat h_l^{(t)}$ in the activation diagnostic, and overload of $N$ for training examples, categories, evaluation prompts, and random directions.  The token steps are discrete and should be written $t\in\{1,\ldots,100\}$ rather than $t\in[1,100]$.  Terminology drifts among ``truthfulness,'' ``factual preference,'' and RefPref, and among ``norm elevation,'' ``saturation,'' ``stabilization,'' and ``manifold departure.''

\section*{Meaningful editorial, compilation, and reference findings}

\subsection*{The source compiles but contains visible Markdown and table overflow}

\findinglabel{Location.} Experimental Setup; Tables~\texttt{tab:long\_gen} and \texttt{tab:mcnemar}; clean pdfLaTeX log.

\findinglabel{Problem.} Experimental Setup still contains Markdown such as \verb|**Dual-Evaluation Paradigm**| and plain numbered lines; the compiled PDF visibly prints the asterisks.  After two pdfLaTeX passes, the manuscript builds to six pages, but the log reports overfull boxes of approximately 61.44~pt for the natural-completion table and 5.54~pt for the McNemar table.

\findinglabel{Why it matters.} The source is technically buildable but not professionally typeset, and the natural table can intrude far into the adjacent IEEE column.

\findinglabel{Suggested fix.} Replace Markdown with LaTeX commands and an \texttt{enumerate} environment.  Redesign or resize both tables while preserving legibility, rebuild twice, and inspect the final log and PDF.

\subsection*{Citation fit remains uneven}

\findinglabel{Location.} Introduction; bibliography entries \texttt{ref10}, \texttt{ref11}, and \texttt{ref12}.

\findinglabel{Problem.} \texttt{ref10} is the LoRA paper rather than direct evidence for broad SFT hallucination mitigation.  \texttt{ref11} concerns irrelevant retrieved context but does not establish that internal activation pathways cause evidence ignoring.  \texttt{ref12} addresses catastrophic forgetting generally rather than the claimed medical-language-model failure mode.

\findinglabel{Why it matters.} Direct evidence and mechanistic hypotheses are presented as if they have the same citation support.

\findinglabel{Suggested fix.} Add direct sources for hallucination-focused SFT, language-model forgetting after domain adaptation, and evidence use in RAG.  Present the activation explanation as a hypothesis unless directly measured or supported.

\section*{Proofing sweep}

The bundled mechanical scan was run on the source and the clean compiled PDF and returned no high-confidence regex hits.  Manual source, equation-adjacent, table, and bibliography checks found the following bounded defects:

\begin{itemize}
    \item \textbf{Contribution~4 and RAG Results:} change ``an non-oracle'' to ``a non-oracle.''
    \item \textbf{Contribution~2 and natural Results:} replace ``un-truncated'' with ``untruncated.''
    \item \textbf{Experimental Setup:} replace Markdown bold and plain numbered text with valid LaTeX.
    \item \textbf{Activation prose:} replace ``within 0.19 units'' and ``$+2.93$ magnitude distortion'' with values computed from the corrected data; specify the hard-cutoff time indices.
    \item \textbf{RAG prose:} distinguish retrieval recall from retrieval miss/noise, and correct the stale latency and context-expansion quantities.
\end{itemize}

The bibliography was checked for duplicate punctuation, malformed quotation punctuation, and obvious citation-format glitches; no additional mechanical defect was found.

\section*{Items requiring external verification}

\begin{itemize}
    \item Verify all formulary-derived reference answers and synthetic counter-answers against the official source, current clinical guidance, qualified clinician review, and applicable reuse permissions.
    \item Validate multilingual BERTScore for Vietnamese clinical negation, quantities, units, contraindications, and interaction mechanisms before interpreting RefPref as factual alignment.
    \item Verify novelty against current token-dependent, schedule-based, and dose-controlled activation-steering work.
    \item Verify the description of Qwen2.5-7B-Instruct as a ``Small Language Model'' and the privacy/local-hospital claims against explicit definitions, a threat model, and measured resource requirements.
    \item Verify BM25 and stronger Vietnamese retrieval baselines before making a general comparison with RAG.
\end{itemize}

\section*{Highest-priority revision sequence}

\begin{enumerate}
    \item Correct and regenerate the activation, latency, prompt-token, corpus-size, and retrieval summary numbers from raw artifacts.
    \item Replace the obsolete natural-completion Discussion and add paired uncertainty for all schedules and metrics.
    \item Redesign the activation analysis so it tests an operational definition of saturation on controlled contexts.
    \item Add clinician-adjudicated correctness/harm endpoints or narrow all claims to reference preference.
    \item Fully specify and strengthen BM25/oracle RAG, random-direction inference, and systems timing.
    \item Add ablation uncertainty, correct the matched-dose sign definition, and test vector-extraction/layer alternatives.
    \item Complete dataset provenance, leakage controls, metric/hook specifications, LaTeX formatting, table layout, and citation fit.
\end{enumerate}

\end{document}